% CS340 Operating Systems -- replication report
% Target paper: Tu et al., "Understanding Real-World Concurrency Bugs in Go", ASPLOS'19.
%
% The report cover comes from docs/cover_template.docx (exported as
% report/cover.pdf). If report/cover.pdf is present it is inserted as page 1;
% if it is missing the document still compiles without the cover.

\documentclass[11pt]{article}

\usepackage[margin=1in]{geometry}
\usepackage{graphicx}
\usepackage{booktabs}
\usepackage{amsmath}
\usepackage[hidelinks]{hyperref}
\usepackage{pdfpages}

\newcommand{\eg}{e.g.}
\newcommand{\ie}{i.e.}
\newcommand{\etal}{et al.}

\title{\bfseries Understanding Real-World Concurrency Bugs in Go:\\
A Replication}
\author{Group 5:\\
辛相天 (1240006592) \\ 张歆海 (1240012291) \\ 谢昊桐 (1230022122) \\ 郭子嘉 (1240015875)}
\date{2026-11-15}

\begin{document}

\IfFileExists{cover.pdf}{\includepdf[pages=1,pagecommand={}]{cover.pdf}}

\maketitle
\thispagestyle{empty}

\begin{abstract}
This report describes our replication of Tu \etal{}\cite{go-study}, the first
empirical study of concurrency bugs in real Go programs. We re-derive the
paper's statistics from the released dataset, re-run a subset of the bugs
against Go's built-in detectors, and extend the study by asking whether
Go~1.22's per-iteration loop-variable semantics removes the bug class the
paper attributes to closures over shared loop variables.

Our level-1 results reproduce the paper's taxonomy. Over the 169
commit-backed bug records in the released workbook we recover 83 blocking and
86 non-blocking bugs and 105 shared-memory and 65 message-passing labels, and
we show that every difference from the paper's published tables is explained
by exactly two documented causes: two issue-only rows that carry no commit
hash, and one record that carries two cause labels where the paper assigns
one. Under the paper's own assignment convention our recomputation matches
paper Tables~5, 6 and~7 cell by cell. Level-2 detector coverage is not yet
measured and is reported as such.
\end{abstract}

\section{Introduction and Claims}

\subsection{Problem}

Go was designed for concurrent programming from the start: goroutines,
channels and a small, uniform set of synchronisation primitives replaced
the ad-hoc mix of threads, mutexes and condition variables that dominated
server-side C and C++. The design rests on an explicit bet that
\emph{message passing is safer than shared-memory synchronisation}, and the
language gives programmers no way to reach into another goroutine's memory
without a channel or an explicit lock. That bet shaped a decade of
production Go code: Docker, Kubernetes, etcd, CockroachDB, gRPC and BoltDB
all rely on it. But the bet is a hypothesis about \emph{how bugs arise},
not a guarantee about how they are prevented, and Go's own idioms introduce
failure modes that neither the C/C++ nor the Erlang/Java literature covers.

\cite{go-study} was the first study to test that hypothesis empirically. It
collected 171 concurrency bugs from six production Go programs, labelled
each one along two independent axes --- \emph{cause} (shared-memory access
versus message passing) and \emph{behaviour} (blocking versus non-blocking)
--- and found that more than half of them are caused by problems specific to
Go rather than by traditional shared-memory races. The practical question we
replicate is therefore narrow and answerable: how much of that picture
survives an independent re-derivation, and does the language's own recent
semantic change close the largest identified bug class?

\subsection{Claims}

The four claims we test are stated, with the paper value and our value, in
\texttt{docs/claims.md}; that file is the contract between the numbers and
the text below, and no number appears in this report unless it appears
there.

\paragraph{C1 (pervasiveness and taxonomy).} Go concurrency bugs are
pervasive in production Go code and require a two-dimensional taxonomy:
171 bugs across six production programs, split 85 blocking / 86 non-blocking
by behaviour and 105 shared-memory / 66 message-passing by cause. Our
commit-backed scope covers 169 of those 171 records; the two excluded rows
and every difference from the published tables are accounted for in
\S\ref{sec:discrepancies}.

\paragraph{C2 (non-traditional causes).} Most of these bugs are not classic
shared-memory races but Go-specific idioms --- misuse of Go libraries and
goroutine/channel patterns --- so message passing is not automatically safer.

\paragraph{C3 (fix strategies).} Fix strategies vary with bug causes: bugs
whose cause involves a particular synchronisation primitive are usually fixed
by adjusting that same primitive. This is a statement about the
\emph{distribution} of fix strategies, and it is what our data supports.

\paragraph{C4 (detection gap).} Go's two built-in concurrency detectors
catch only a minority of the bugs they are aimed at. The paper reports the
deadlock detector firing on 2 of 21 reproduced blocking bugs, and the race
detector on 10 of 20 reproduced non-blocking bugs within 100 runs each ---
exactly half of that sample. Our own coverage is not yet measured.

\subsection{Scope}

Our reproduction has two layers, mirroring the structure of the paper.

\paragraph{L1 --- re-derive the statistics.} We parse the released dataset
at a pinned upstream commit and recompute every table ourselves, then
compare against the published numbers: the behaviour/cause taxonomy and
application breakdown (paper Tables~5, 6, 7), the root-cause and fix-strategy
statistics for non-blocking bugs (Tables~9, 10). We do not copy the paper's
rendered tables; the taxonomy and counting logic are re-implemented from the
paper's definitions and applied to the released workbook.

\paragraph{L2 --- re-run the bugs.} We rebuild a subset of 15--20 selected
bugs at their buggy commits, trigger each one, and re-measure what Go's
built-in deadlock and race detectors report. This covers the paper's bug
reproduction and detector evaluation (Tables~8 and~12, sections~5.3 and
6.3). Level~2 has not been run yet, so this report presents our level-1
results in full and states the level-2 gap explicitly rather than filling it
with estimates.

\paragraph{Extension.} We change exactly one factor: the Go language
version. The same bug is rebuilt under Go~1.21 and Go~1.22 to test whether
the per-iteration loop-variable semantics introduced in Go~1.22 eliminate the
anonymous-function bug class the paper singles out.

\section{Reproduction Setup}
\label{sec:setup}

\subsection{Dataset}

All level-1 numbers come from one pinned artifact.

\begin{table}[h]
\centering
\caption{The pinned dataset and its identity.}
\begin{tabular}{@{}ll@{}}
\toprule
Item & Value \\
\midrule
Artifact & \texttt{bug table.xlsx} \\
Upstream repository & \texttt{system-pclub/go-concurrency-bugs} \\
Pinned commit & \texttt{f2655063a2278cff64a51121b73801eda05fdcfa} \\
SHA-256 & \texttt{a0e5f0103600b3c394b86cb975f403789366f0190312019fae648de0cd8e3235} \\
Local copy & \texttt{data/bug\_table.xlsx}, unmodified \\
Entries in the workbook & 171 \\
Records in our inventory & 169 (83 blocking, 86 non-blocking) \\
\bottomrule
\end{tabular}
\end{table}

The fetcher verifies the workbook's checksum and never silently replaces a
mismatched file; the parser runs offline and re-verifies the same checksum.
The parsing is implemented from scratch with the Python standard library only
--- no spreadsheet library, no third-party package --- so the counts do not
depend on how a given tool interprets a worksheet. Column positions are
declared explicitly rather than inferred, because the two worksheets have
different layouts, and header cells are checked before any cell is read so
that a layout change fails loudly instead of silently shifting a column.

The parser is exposed through the repository's \texttt{make} targets:

\begin{center}
\texttt{make setup}\quad \texttt{make check}\quad \texttt{make tables}
\end{center}

and its behaviour is covered by a regression test suite that exercises
corrupt and missing workbooks, blank flags, overlapping labels, sparse and
inline strings, and stable byte-for-byte CSV output.

\subsection{Host and toolchain}

Level-1 statistics depend only on the Python standard library and are
therefore identical on every member's machine; \texttt{make tables} completes
in under two seconds. The level-2 experiments are a different matter: hang
detection, timeout behaviour, signal delivery and the race detector are
operating-system dependent, so they must run on Linux.

\begin{table}[h]
\centering
\caption{Execution environment. The level-2 rows are \emph{pending}: no image
has been built and verified yet, and no level-2 number appears in this
report until they are filled in.}
\begin{tabular}{@{}ll@{}}
\toprule
Item & Value \\
\midrule
Level-1 host & Windows 11, Intel Core 5 220H, 16 logical cores, 15.6\,GB \\
Level-1 runtime & Python 3.11.15, standard library only \\
Experiment image & \texttt{golang:1.22.4-bookworm} (Go~1.22) \\
Extension image & \texttt{golang:1.21.13-bookworm} (Go~1.21) \\
Base image digest & \emph{pending --- first build} \\
Kernel (\texttt{uname -a}) & \emph{pending --- first build} \\
Compiler (\texttt{gcc --version}) & \emph{pending --- first build} \\
\texttt{GOMAXPROCS} & 4, fixed in the \texttt{Dockerfile} \\
\texttt{CGO\_ENABLED} & 1 (the race detector requires cgo) \\
\texttt{GOTOOLCHAIN} & \texttt{local}, so the image's own toolchain is used \\
Detector flags & \texttt{-race} for data races; the deadlock detector is always on \\
Runs per bug & $\geq 3$ for statistics; 100 for the race-detector experiment \\
\bottomrule
\end{tabular}
\end{table}

The two images differ only in the Go version, which is the single variable in
our extension experiment. The base image must ultimately be pinned by digest
rather than by tag, so that an upstream retag cannot silently change what we
ran.

\subsection{Protocol}

Each reproduced bug is rebuilt at the commit that the dataset records as its
fixing commit's first parent, which is the pre-fix revision. Each candidate
carries a provenance record naming the repository, the full pre-fix and
fixing commit hashes, the inventory row it came from, the exact trigger, the
oracle used to decide whether the symptom occurred, and the known limits of
the reproduction. The upstream checkout itself is not committed, and neither
are binaries or raw logs; only the specification, the harness and a summary
are tracked.

A reproduction is only accepted once it has produced a build log and three
runs whose detector outcome is recorded; the Go version, the image digest,
\texttt{GOMAXPROCS}, \texttt{CGO\_ENABLED} and the run count travel with the
number. An unlabelled run log is not a result.

\section{Results}
\label{sec:results}

\subsection{Level 1: the taxonomy reproduces}

Table~\ref{tab:t5} compares our recomputed taxonomy against the paper's
published Table~5. The paper's columns are the behaviour split and the cause
split; ours are computed from the pinned workbook.

\begin{table}[h]
\centering
\caption{Behaviour and cause taxonomy, per application. The published
Table~5 is shown beside our recomputation over the 169 commit-backed
records. Differences are confined to BoltDB and Kubernetes and are explained
in \S\ref{sec:discrepancies}.}
\label{tab:t5}
\begin{tabular}{@{}lrrrrrrrr@{}}
\toprule
& \multicolumn{2}{c}{blocking} & \multicolumn{2}{c}{non-blocking}
& \multicolumn{2}{c}{shared memory} & \multicolumn{2}{c}{message passing} \\
\cmidrule(lr){2-3}\cmidrule(lr){4-5}\cmidrule(lr){6-7}\cmidrule(lr){8-9}
Application & Paper & Ours & Paper & Ours & Paper & Ours & Paper & Ours \\
\midrule
Docker      & 21 & 21 & 23 & 23 & 28 & 28 & 16 & 16 \\
Kubernetes  & 17 & 17 & 17 & 17 & 20 & \textbf{21} & 14 & 14 \\
etcd        & 21 & 21 & 16 & 16 & 18 & 18 & 19 & 19 \\
CockroachDB & 12 & 12 & 16 & 16 & 23 & 23 &  5 &  5 \\
gRPC        & 11 & 11 & 12 & 12 & 12 & 12 & 11 & 11 \\
BoltDB      & \textbf{3} & \textbf{1} & 2 & 2 & \textbf{4} & \textbf{3} & \textbf{1} & \textbf{0} \\
\midrule
Total       & \textbf{85} & \textbf{83} & 86 & 86 & 105 & 105 & \textbf{66} & \textbf{65} \\
\bottomrule
\end{tabular}
\end{table}

Two observations matter. First, the shared-memory total is unchanged at 105
even though two individual rows differ, because the two changes run in
opposite directions and cancel at the total; a total-only comparison would
have hidden both. Second, the four cells that differ are exactly the cells
that the two causes in \S\ref{sec:discrepancies} predict, which is a
stronger statement than agreement in aggregate.

Table~\ref{tab:t6} gives the causes of blocking bugs, our recomputation of
paper Table~6.

\begin{table}[h]
\centering
\caption{Causes of blocking bugs, per application, beside the published
Table~6.}
\label{tab:t6}
\begin{tabular}{@{}lrrrrrrrr@{}}
\toprule
& \multicolumn{3}{c}{Shared memory} & \multicolumn{3}{c}{Message passing}
& \multicolumn{2}{c}{Total} \\
\cmidrule(lr){2-4}\cmidrule(lr){5-7}\cmidrule(lr){8-9}
Application & Mutex & RWMutex & Wait & Chan & Chan w/ & Lib & Paper & Ours \\
\midrule
Docker      & 9 & 0 & 3 & 5 & 2 & 2 & 21 & 21 \\
Kubernetes  & 6 & 2 & \textbf{1} & 3 & 6 & 0 & 17 & 17 \\
etcd        & 5 & 0 & 0 & 10 & 5 & 1 & 21 & 21 \\
CockroachDB & 4 & 3 & 0 & 5 & 0 & 0 & 12 & 12 \\
gRPC        & 2 & 0 & 0 & 6 & 2 & 1 & 11 & 11 \\
BoltDB      & \textbf{1} & 0 & 0 & 0 & 0 & 0 & \textbf{3} & \textbf{1} \\
\midrule
Total       & \textbf{27} & 5 & \textbf{4} & 29 & \textbf{15} & 4 & \textbf{85} & \textbf{83} \\
\bottomrule
\end{tabular}
\end{table}

Our \texttt{Wait} column is one higher and our \texttt{Chan w/} column one
lower than the paper's, for the reason given in \S\ref{sec:discrepancies}: one
record carries both a \texttt{Wait} and a \texttt{Chan w/} cause, and the
paper assigns it only one of them. Note also that our cause column totals
84 assignments over 83 blocking bugs; the totals are category assignments,
not deduplicated bug counts.

For the non-blocking half of the study we reproduce the published Table~9
(root causes) in \emph{all 42 cells} and the published Table~10 (fix
strategies) in \emph{all 35 cells}, including both total rows. This is the
part of the paper's dataset that carries no overlapping labels, and it
reproduces without qualification.

\subsection{Level 1: non-traditional causes}

The paper's central qualitative finding is that more than half of the bugs
are not traditional shared-memory races but problems specific to Go. The
non-traditional label is a \emph{separate} classification from the
shared-memory / message-passing cause axis of Table~\ref{tab:t5}, and the
two must not be substituted for one another: the share of bugs caused by
shared-memory access is not the share of bugs that are non-traditional.

Counting the workbook's own non-traditional label (column~L for blocking
rows, column~H for non-blocking rows) over our 169 commit-backed records
gives 54 of 83 blocking rows and 40 of 86 non-blocking rows, that is
\textbf{94 of 169, or 55.6\,\%}, carrying a non-traditional label. This
supports a majority claim \emph{within this dataset}. It does not establish
a universal failure probability for Go constructs, and it is a supplementary
workbook-label count rather than a column of the five published tables; the
source coordinates travel with the number wherever it is used.

\subsection{Level 2: not yet measured}

The paper reports that Go's built-in deadlock detector fires on 2 of 21
reproduced blocking bugs (Table~8), and that the race detector fires on 10
of 20 reproduced non-blocking bugs when each is run 100 times (Table~12).
The race detector therefore detects exactly half of that sample; it is not
accurate to say that it misses more than half of it.

We report no detector coverage of our own. The candidate selection and the
reproducer contract are complete, but no bug has been built and triggered
yet, because the experiment image has not been built and verified on a
machine with Docker. Until that is done, quoting a coverage figure would be
an estimate presented as a measurement.

One discrepancy is already visible in the upstream annotations and is
recorded now rather than later: the commit-backed inventory marks 19
blocking and 23 non-blocking records as reproduced, whereas the paper reports
21 and 20. Including the two excluded issue-only rows makes the workbook's
blocking flag count 21, matching the paper, so the blocking difference is a
scope effect; the 23 non-blocking flags against the paper's 20 remain
unexplained. These flags are upstream annotations, not evidence that this team
has reproduced anything.

\section{Analysis of Discrepancies}
\label{sec:discrepancies}

Every difference between our numbers and the paper's has one of two causes.
Both are properties of the released workbook, not artefacts of our counting.

\subsection{Two issue-only rows carry no commit}

The workbook holds 171 bug entries. Two blocking rows have an issue ID where
a commit hash should be:

\begin{center}
\begin{tabular}{@{}llll@{}}
\toprule
Sheet & Row & Issue ID & Cause \\
\midrule
blocking & 90 & \texttt{boltdb\#392} & \texttt{Mutex}, shared memory \\
blocking & 92 & \texttt{boltdb\#240} & \texttt{Chan w/}, message passing \\
\bottomrule
\end{tabular}
\end{center}

These rows are \emph{present} in the released workbook. They are not missing
data and not an upstream omission; they are the two entries the artifact does
not bind to a commit. Because we can only reproduce a bug we can check out,
our inventory excludes them, which accounts for $-2$ blocking, $-1$
shared-memory and $-1$ message-passing relative to the paper.

That exclusion is correct for our purpose, and it is also independently
corroborated: these two issue IDs are exactly the two bugs that the paper's
built-in deadlock detector reports. The paper did analyse them --- it simply
did so through issue identifiers rather than through commits.

\subsection{One record carries two cause labels}

One included record, \texttt{kubernetes\#11298}, carries \emph{both} a
shared-memory and a message-passing label, and \emph{both} a \texttt{Wait}
and a \texttt{Chan w/} cause, in the released workbook. The paper assigns it
a single cause: message passing. We retain both labels, because discarding a
label that the dataset contains would be a silent correction of the source
data, and the parser warns about the overlap rather than hiding it.

This single record explains the remaining differences: $+1$ shared-memory
label for Kubernetes, $+1$ \texttt{Wait} cause, and the corresponding $-1$ in
the \texttt{Chan w/} column of the blocking-cause table and in the
\texttt{Chan w/} / \texttt{Remove} cell of the fix-strategy table.

\subsection{The two causes account for every difference}

The point of attributing differences individually is that they can be checked
against the paper's own tables rather than merely asserted. Under the
paper's assignment convention --- include the two issue-only rows, and give
\texttt{kubernetes\#11298} one cause --- our recomputation reproduces
paper Table~5 cell by cell for all six applications, paper Table~6 in all
six cause columns, and paper Table~7 in every cell. Tables~9 and~10 match
under either convention.

We still report the commit-backed numbers as our primary result, because
those are the ones backed by a commit we can actually check out and rebuild.
The reconciliation is a consistency argument, not a licence to present
counts we cannot reproduce.

\subsection{Setup deviations}

Beyond the dataset, the report must record the setting differences between
the paper and our study.

\begin{table}[h]
\centering
\caption{Setup deviations from the paper.}
\begin{tabular}{@{}p{0.16\textwidth}p{0.22\textwidth}p{0.24\textwidth}p{0.24\textwidth}@{}}
\toprule
Item & Paper setup & Our setup & Reason \\
\midrule
Go version & older toolchain, 2015--2018 code & Go 1.22.4 & one unified
  environment; old code may need API repairs, recorded per bug \\
Reproduction scale & 41 bugs re-run & subset, target 15--20 & limited time
  and manpower; the level-1 counting metrics are insensitive to scale \\
Hardware & the authors' machines & Linux container on a laptop & the metrics
  are counts, not throughputs; stated as a limitation \\
Detector protocol & 1 run per blocking bug, 100 per non-blocking bug & same,
  plus $\geq 3$ runs for stability & our extra stability requirement is
  additional to the paper's protocol \\
\bottomrule
\end{tabular}
\end{table}

\subsection{Claims we deliberately do not make}

Two claims that would have been easy to write are withheld because the
current work does not support them.

First, the paper's fix data does not by itself establish that upstream fixes
are shallow, incomplete, or that they introduce new bugs. Counting small
patches or added synchronisation cannot show that developers failed to
identify a root cause. Our level-1 data supports the narrower and
well-defined claim that fix strategies vary with bug causes, and a
before/after reproducer control can later test whether a specific fix
removes a specific symptom --- but a per-bug control cannot prove a general
claim about fix quality.

Second, the paper's non-blocking race-detector result is a detection rate of
exactly half, not a failure rate above half. Reporting it the other way
round would overstate the paper's result in the paper's favour and against
Go's tooling.

\section{Extension}
\label{sec:extension}

\subsection{Hypothesis}

The paper observes that some non-blocking bugs are caused by anonymous
functions that capture a loop variable shared across iterations. Go~1.22
changed the semantics of \texttt{for} loops: each iteration now receives its
own variables, subject to the language version selected for the module
\cite{go122,loopvar}. That change postdates the 2019 paper, so it is not a
claim the paper makes and not a result we may claim until we have run it.

\subsection{Protocol}

We change exactly one factor, the Go language version, and hold everything
else fixed: the same source revision, the same trigger, the same workload,
the same \texttt{GOMAXPROCS} and \texttt{CGO\_ENABLED} settings, and the same
detector invocation. The only difference between the two runs is the image
tag, \texttt{golang:1.21.13-bookworm} against \texttt{golang:1.22.4-bookworm}.

Three controls must be recorded with every run, because changing the
container tag alone does not establish that the intended language semantics
changed:

\begin{enumerate}
\item the module's \texttt{go} directive and the actual toolchain version,
  since the language version in effect follows the directive and not the
  image tag;
\item any compiler flags used in the build;
\item the base image and compiler differences, recorded rather than assumed.
\end{enumerate}

\subsection{Scope of the hypothesis}

The hypothesis is narrow on purpose: the new loop-variable semantics should
remove selected \emph{loop-capture} symptoms. Anonymous-function bugs that
involve other shared state --- a field, or a closure over a variable
declared outside the loop --- need not disappear at all, and we will report
them as surviving if they do. A result in which only the loop-capture subset
is fixed is a useful result, and reporting it precisely is more informative
than claiming the bug class is closed.

\section{Conclusion}

We set out to re-derive the statistics of the first empirical study of Go
concurrency bugs and to ask whether Go~1.22 closes the loop-variable bug
class. On the first half, the replication succeeds. Over the 169
commit-backed records of the pinned dataset we recover 83 blocking and 86
non-blocking bugs and 105 shared-memory and 65 message-passing labels, we
reproduce paper Tables~9 and~10 in every cell, and we account for every
remaining difference with two documented properties of the released
workbook. The 94 of 169 records carrying a non-traditional label, 55.6\,\%,
independently support the paper's central qualitative finding that Go's own
idioms, not classical shared-memory races, are the majority cause of Go
concurrency bugs --- and that message passing is therefore not the safer
choice it is often assumed to be.

Two things are worth stating plainly rather than leaving to be inferred.
First, the numbers in this report are only as good as their provenance, and
the single most useful thing we did was refuse to report a total without
asking which cells moved: the shared-memory total agrees with the paper
exactly, and that agreement conceals two offsetting row-level differences.
Second, some attractive claims are simply not supported by the evidence
available to us. That upstream fixes are shallow or that they introduce new
bugs is not something a fix-strategy table can show, and the paper's race
detector detected exactly half of its sample rather than missing more than
half.

The second half of the programme --- rebuilding the bugs, measuring what
Go's built-in detectors actually report, and testing the loop-variable
hypothesis --- is not yet done. The candidate selection, the reproducer
contract and the environment specification are in place, but no bug has been
built and triggered and the experiment image is not yet built and verified.
We therefore report no detector coverage of our own, and the extension is
stated as a protocol rather than as a finding. The next step is the
environment hand-off: build both images on a Linux host, pin the base image
by digest, record the kernel and compiler versions, and only then produce
the first reproducer logs. Everything downstream --- our own coverage
numbers, the detector comparison, and the Go~1.22 result --- becomes
quotable at that point and not before.

\section*{Contribution Statement}

All four members contributed to this work, and the per-commit breakdown is
recorded in \texttt{docs/weekly-log.md}, which lists, for every commit, the
author, the date and the report section that commit produced material for.

辛相天 implemented the pinned dataset pipeline and re-derived Tables
5/6/7/9/10; 张歆海 built and pinned the execution environment; 谢昊桐
rebuilt and re-ran the selected bugs for the L2 experiments; 郭子嘉
maintained the claims file, this report, and the contribution log. The
individual extent of each contribution is verifiable from the git history:
every commit carries the author's student ID and full name in its subject.

Where a member's work spans roles, the weekly log records who wrote the
code rather than only who owns the analysis: the pinned-data fetcher and
offline parser were implemented by 谢昊桐 as a cross-role prerequisite at
the team's request, while the statistical scope and the deviations register
remain 辛相天's.

\section*{AI Usage Statement}

We used AI coding assistants while preparing this repository, and we
disclose them here as a group so that the statement matches what the
members reported in their own records.

An OpenClaw coding agent assisted with the documentation and report work
under 郭子嘉's direction: the role table, the contribution log, the claims
file, and the prose and structure of this report. Codex assisted 谢昊桐 with
the data scripts and their tests, the inspection of upstream source, the
candidate selection, the reproducer specification, and the drafting of the
review documents.

The use of these tools was limited to writing and editing text and code.
Every technical content in this report was verified by the human members
before being committed: each number was produced by the code in this
repository --- \texttt{make tables} for level~1, and \texttt{repro/} run logs
for level~2 --- and re-checked against the paper. No level-2 number appears
in this report, precisely because no run log exists yet. The AI assistants
did not run the experiments, did not choose the reproduced bugs, and did not
make claims about the paper that are not traceable to
\texttt{docs/claims.md}. The claims, the deviations and the conclusions are
the members', and the git history attributes them.

\bibliographystyle{plain}
\bibliography{refs}

\end{document}
